\section{Experimental Setup}
\label{sec:setup}

\paragraph{Models and decoding.}
The main reader is Qwen2.5-7B-Instruct~\cite{qwen25}; the cross-model check uses
Mistral-7B-Instruct-v0.3~\cite{mistral7b}. Both run in fp16 with greedy decoding,
$64$ new tokens, maximum context $24{,}576$ tokens and a strict policy that aborts rather
than truncates when a prompt exceeds it. The random seed is fixed at \texttt{20261015}.
Scoring uses the official SQuAD~v2 script, a fixed, model-independent procedure.

\paragraph{Baselines and fairness.}
We re-run six published compressors through their reference implementations:
LLMLingua-2~\cite{llmlingua2}, LongLLMLingua~\cite{longllmlingua}, Selective
Context~\cite{selective_context}, Provence~\cite{provence}, CPC~\cite{cpc}, and
Quito-X~\cite{quitox}. We also run a 2026 rule-based compressor that needs no model at
compression time~\cite{chokshi-lexical-2026}, the release that accompanies the
peer-reviewed study of linguistic rules as compressors~\cite{ma-linguistic-2026};
it is a preprint arm and is labelled as such. Where a method exposes both a
budget-matched and a native configuration we report both; the budget-matched one is used for the headline comparison,
since published defaults vary in length. We additionally run a masking wrapper around
LLMLingua-2, and our own non-routing arms: \arm{full},
\arm{alias} (alias all protected values, no compression), \arm{delete} (redact all, no
compression), \arm{alias\_budget} and \arm{alias\_budget\_link} (global alias, then
redundancy-based selection without and with anchor protection), and the corresponding
\arm{budget\_cluster} pair on raw text.

All compressors receive the question and the context and nothing else. We verified in the
runner that the reference answer is never passed to any baseline, and that no baseline
consumes the anchor set. This is the fairness condition for the comparison: the
contribution we claim is precisely the ability to obtain the referential benefit without
the answer key.

\paragraph{Budgets and statistics.}
Budget arms are anchored at $r=0.5$ for the main comparison and swept over
$r\in\{0.10,0.15,0.25,0.50,0.75\}$ for the budget curve. Published compressors are
evaluated at their own $50\%$ setting, with the token accounting reported in every table so
that a method that realises a different length is visible as such.

All comparisons between two arms on the same cases are paired. We report the mean paired
difference and a $95\%$ percentile bootstrap interval with $10{,}000$ resamples over cases.
A difference is called significant when the interval excludes zero.

Three verdicts are kept apart throughout, because they license different sentences. A
difference is \emph{larger} or \emph{smaller} when its interval excludes zero. It is
\emph{not significantly different} when the interval contains zero---which is the verdict for
\textsc{DARK} against Quito-X on F1, and it supports no more than the statement that this
experiment did not distinguish them. We do not claim non-inferiority anywhere: that would
require a margin fixed in advance and a one-sided test against it, and neither is present here.
Where a result is reported as a win it is a difference whose interval clears zero, and we say so
where it is first used. We deliberately do not
attach a significance claim to any comparison across the two benchmark sizes, which draw
different cases.

\paragraph{Reproducibility controls.}
Two controls are run in every batch. First, the uncompressed arm is recomputed; it involves
no selection logic, so it must reproduce the reference value, and any deviation indicates
that the run differs in model, precision, or data. Second, every budget arm asserts
non-overshoot at the case level. Runs that fail either control are discarded rather than
reported.
